% ==============================================================================
%  The learned operator, audited -- standalone report
%
%  Companion to reports/glueball/glueball_spectroscopy.tex and
%  reports/attention/attention_as_operator.tex; self-contained.
%  Build: pdflatex (run twice for refs).
%  Figures are PNGs under ../../results/<study>/; \graphicspath lists those dirs.
% ==============================================================================
\documentclass[11pt,a4paper]{article}

\usepackage[margin=2.6cm]{geometry}
\usepackage[T1]{fontenc}
\usepackage[utf8]{inputenc}
\usepackage{lmodern}
\usepackage{microtype}
\usepackage{amsmath,amssymb,mathtools}
\usepackage{graphicx}
\usepackage{booktabs}
\usepackage{caption}
\usepackage{xcolor}
\usepackage[colorlinks=true,linkcolor=blue!50!black,citecolor=blue!50!black,urlcolor=blue!50!black]{hyperref}

\emergencystretch=3em

% ---- float placement: this report is table-heavy ------------------------------
\renewcommand{\topfraction}{0.92}
\renewcommand{\bottomfraction}{0.75}
\renewcommand{\textfraction}{0.07}
\renewcommand{\floatpagefraction}{0.80}
\setcounter{topnumber}{3}
\setcounter{totalnumber}{4}

\graphicspath{{../../}{../../results/}{../../results/sampler/}{../../results/glueball/}{../../results/attention/}{../../results/fair_fight/}{../../results/wilson_regression/}}

% ---- shorthand ---------------------------------------------------------------
\newcommand{\Tr}{\operatorname{Tr}}
\newcommand{\Obar}{\bar{O}}
\newcommand{\meff}{m_{\mathrm{eff}}}
\newcommand{\at}{a_t}
\newcommand{\as}{a_s}
\newcommand{\That}{\hat{T}}
\newcommand{\Azero}{A_0}
\newcommand{\dAzero}{\Delta A_0}
\newcommand{\OG}{O_{\mathrm{GELT}}}
\newcommand{\spn}{\operatorname{span}}
\newcommand{\arm}[1]{\texttt{#1}}

\title{One Learned Operator, Three Audits:\\
the Input--Architecture Curve, the Strengthened Comparator,\\
and What the Network Actually Found}
\author{Francesco Passante}
\date{September 12, 2026}

\begin{document}
\maketitle

\begin{abstract}
The $SU(2)$ spectroscopy chapter of this thesis rests on one number: a
gauge-equivariant lattice transformer trained as a variational $0^{++}$
operator carries more ground-state weight than the optimal combination of the
classical smeared basis, $\dAzero = +0.078 \pm 0.022$ ($3.6\sigma$) over two
independently sampled ensembles at a consistent mass. A single number of that
shape invites three separate objections, and this report is the record of
answering all three offline, on the same $400$ held-out configurations per
ensemble, with the estimator and the readings fixed in writing before the
numbers existed. \textbf{Was the opponent strong enough?} Against the
input-matched seven-level classical ladder the claim reads
$+0.095 \pm 0.021$; against the strongest basis the repository builds --- $21$
operators, seven smearing levels times three loop shapes --- it falls to
$+0.038 \pm 0.027$, so genuine loop-shape variety closes about $60\%$ of the
gap and leaves a $1$--$2\sigma$ edge. Reading that arm at all required dropping
the near-null directions of $C(t_0)$ rather than flooring them, and the
pre-registered rule for choosing between the two turned out to have no answer,
for a reason worth recording: the GEVP at $t_0=1$ maximises $C(2)/C(1)$ while
the variational gate tests $C(2)/C(0)$. \textbf{Is it the architecture or just
richer inputs?} Replacing the single number by $\Azero$ as a function of the
input content handed to \emph{both} methods gives a curve on which the learned
operator sits above the classical one at every resolvable point
($+0.252 \pm 0.056$ on thin links, $+0.075 \pm 0.019$ at four smearing levels),
an untrained network of identical architecture sits below the classical one
everywhere, and a width control shows the gain is input content and not
capacity. None of the three pre-registered readings survives its own test
unmodified. \textbf{What did the network find?} Because the equal-time
connected correlator \emph{is} the Hilbert-space inner product, the split
$\OG = P + r$ into a classical projection and an orthogonal remainder is exact.
The same trained operator is $21.1\%$, $12.8\%$ and $2.5\%$ outside the $4$-,
$7$- and $21$-operator spans, deleting $r$ costs the entire advantage
($+0.097 \pm 0.021$ against the strong arm), and $80\%$ of $r$ is rectangular
Wilson loops the network rediscovered. The decisive control is an untrained
network: it is \emph{further} outside the span ($16$--$74\%$) than the trained
one, so the norm fraction is architectural --- what is learned is the amplitude
share $Z_r/Z_G$ ($0.146$ against $0.047$), the fact that the residual becomes
loop-shaped at all ($80\%$ against $0.3\%$), and the sign of
$\dAzero(\text{net}-P)$, which flips to $-0.091 \pm 0.042$ without training.
\end{abstract}

\tableofcontents

% ==============================================================================
\section{The claim, and the three questions it invites}
\label{sec:intro}
% ==============================================================================

The result being audited is stated in
\texttt{reports/glueball/glueball\_spectroscopy.tex} and reproduced here in one
sentence. On an anisotropic $SU(2)$ lattice ($12^3\times 24$, $\beta = 2.4$,
$\xi = 3.0$), a GELT network trained per timeslice on the Rayleigh-quotient loss
$-C(1)/C(0)$ produces a single scalar operator whose ground-state overlap
exceeds that of the variationally optimal combination of a four-level APE
smearing basis:
\begin{equation}
  \dAzero \;=\; \Azero(\text{GELT}) - \Azero(\text{GEVP})
  \;=\; +0.078 \pm 0.022 \quad (3.6\sigma),
  \label{eq:published}
\end{equation}
combining two independently sampled ensembles, at a fitted mass consistent with
unchanged.\footnote{\texttt{notes/glueball\_spectroscopy.md}, ``Replication
results (2026-07-05)''. The two inputs are $+0.066 \pm 0.031$ on the Run-5
ensemble and $+0.089 \pm 0.030$ on the independent seed-1 ensemble.}

Equation~\eqref{eq:published} is a comparison, and every comparison of this
shape admits exactly three lines of attack. They are logically independent, they
require different evidence, and this report devotes one section to each.

\begin{description}
  \item[\textnormal{\textbf{Was the opponent strong enough?}}] The classical arm
        of Eq.~\eqref{eq:published} is four smearing levels of a \emph{single}
        loop shape, the $1\times1$ spatial plaquette. A practitioner building a
        Morningstar--Peardon basis would use many shapes. If a stronger
        classical basis closes the gap, the result is about a weak comparator
        rather than about a network. Section~\ref{sec:fairfight}.
  \item[\textnormal{\textbf{Is it the architecture, or just richer inputs?}}]
        The network is fed smeared link configurations, and one that merely
        \emph{relays} its input advantage is not an interesting result. The
        replacement for a single $\dAzero$ is a \emph{curve}: $\Azero$ as a
        function of the input content handed to both methods, with a third trace
        --- the same architecture, untrained --- separating architecture from
        training. Section~\ref{sec:curve}.
  \item[\textnormal{\textbf{What did the network actually find?}}] Granting the
        win, a referee asks whether the advantage is \emph{new operator content}
        or the \emph{same content in a combination the GEVP failed to locate}.
        The first says the classical variational space has a ceiling; the second
        says the estimator was noisy and more configurations would fix it. The
        question has an exact answer, because operators acting on the vacuum are
        vectors in a Hilbert space whose inner product the correlator already
        computes. Section~\ref{sec:decomp}.
\end{description}

All three audits are pure re-analyses of arrays that already exist: the
test-split time series $\Obar(t)$ of every operator, learned and classical, on
the $400$ held-out configurations of each ensemble. Every number here was
produced on a laptop CPU in seconds to minutes from files tracked in
\texttt{dumps/} --- no GPU time, no training, and, importantly for the first two
audits, no possibility of quietly re-training until the answer improved.

\paragraph{The pre-registration.} Sections~1 and~2 of
\texttt{notes/fable5.1\_10-09\_audit.md} fix the observable, the estimator, the
fit conventions, the grid of measurement points, and three named readings with
their operational tests, and were committed \emph{before} any of the numbers
below existed. This is the only defence against the objection that a curve with
five points and three traces offers many places to stop looking. It is also why
this report says twice, in Sections~\ref{sec:selection} and~\ref{sec:readings},
that a pre-registered rule failed on contact with the data --- recording the
amendment and its reason is the cost of the discipline, and the alternative
(silently choosing what works) is exactly what the discipline exists to
prevent.

% ==============================================================================
\section{Method}
\label{sec:method}
% ==============================================================================

This section is the minimum needed to read Sections~\ref{sec:fairfight}
to~\ref{sec:decomp}. The path integral, the transfer matrix, the spectral
decomposition of a Euclidean correlator, APE smearing, the anisotropic action,
and the heat-bath/overrelaxation sampler are developed from first
principles in the companion report \texttt{glueball\_spectroscopy.tex} and are not
repeated.

\subsection{Correlators, the cosh fit, and the overlap fraction}
\label{sec:a0}

Let $\Obar(t)$ be the zero-momentum, timeslice-summed value of a scalar
operator on one gauge configuration. The $0^{++}$ channel has a nonzero vacuum
expectation value, so the object with a spectral decomposition is the
\emph{connected} correlator
\begin{equation}
  C(\Delta) \;=\; \big\langle\, d(t)\, d(t+\Delta) \,\big\rangle,
  \qquad d(t) = \Obar(t) - \langle \Obar \rangle,
  \label{eq:conn}
\end{equation}
averaged over $t$ and over configurations. On a periodic lattice of temporal
extent $N_t$ the transfer matrix gives
\begin{equation}
  C(\Delta) \;=\; \sum_{n\ge1} |Z_n|^2
  \Big[ e^{-m_n \Delta} + e^{-m_n (N_t-\Delta)} \Big],
  \qquad Z_n = \langle 0 | \Obar | n \rangle .
  \label{eq:spectral}
\end{equation}
Fitting the single-state form $C(\Delta) \approx A\,[e^{-m\Delta} +
e^{-m(N_t-\Delta)}]$ over a window in which the excited states have died gives
$(m, A)$; the mass $m$ is the physics and $A = |Z_0|^2$ is the operator's
coupling to the ground state. The figure of merit for an \emph{operator}, as
opposed to a mass, is the fraction of the operator's total norm that sits on the
ground state,
\begin{equation}
  \boxed{\;\Azero \;=\; \frac{A\,\big(1 + e^{-m N_t}\big)}{C(0)}\;}
  \label{eq:a0}
\end{equation}
--- the numerator is the ground-state contribution to $C(0)$ evaluated from the
fit, the denominator the measured $C(0)$, which by Eq.~\eqref{eq:spectral}
contains every state. A pure ground-state interpolator has $\Azero = 1$;
excited-state contamination pushes it down.

Three properties of $\Azero$ matter for what follows.

\begin{enumerate}
  \item \textbf{It is the right figure of merit at fixed mass.} Two operators
        that agree on $m$ but differ on $\Azero$ measure the same physics with
        different amounts of contamination --- the comparison this report needs,
        and Morningstar and Peardon's own operator-quality metric. It is also
        \emph{exactly what the training loss maximises}: the Rayleigh quotient
        $-C(1)/C(0)$ is minimised by the operator of largest ground-state
        fraction at the mass gap, so objective and metric coincide.
  \item \textbf{$\Azero > 1$ can occur, and is a fit artifact rather than a
        probability violation.} $\Azero$ is a ratio of a \emph{fitted} amplitude
        to a \emph{measured} $C(0)$, estimated on a finite sample; when the
        operator is nearly pure, the two estimates differ by noise in both
        directions. The seed-1 ensemble reads $\Azero = 1.013 \pm 0.062$ for the
        four-level learned operator and $1.066 \pm 0.059$ for the seven-level
        one. These are quoted unclipped, with $1-\Azero$ carrying its sign;
        clipping at one would bias every difference in which they appear.
  \item \textbf{It saturates.} $\Azero \le 1$, and every strong arm on these
        ensembles already reads $0.90$--$1.07$. Two methods separated by less
        than the errors near the ceiling are not thereby shown to be equivalent;
        the observable has simply run out of range. This is pre-registered as
        the \emph{saturation clause} and is enforced mechanically in
        Section~\ref{sec:readings}.
\end{enumerate}

\paragraph{Fit and error conventions.} $(m,A)$ come from
\texttt{fit\_cosh\_correlator}: a profiled-$A$ grid scan with parabolic refine,
dependency-free so that it is safe to call inside every jackknife replica, with
diagonal $\chi^2$ weights $\sigma_\Delta$ fixed from the full sample. Errors are
delete-block jackknife over configurations with block $10$, and the
\emph{entire} fit --- including, for the classical arm, the re-solution of the
GEVP eigenvector --- is redone inside every replica. Differences between two
operators are always \emph{correlated} differences on shared configurations:
the difference is formed inside each replica and jackknifed, not assembled from
two independent error bars.

\subsection{The classical comparator: the GEVP ground vector}
\label{sec:gevp}

A learned operator is one operator. Comparing it against a \emph{basis} would be
unfair in the network's favour; comparing it against the best \emph{single}
member of that basis would be unfair against it. The apples-to-apples object is
the optimal single operator inside the basis's span, and the variational method
supplies it. Given the correlator matrix $C_{ab}(\Delta)$ of a basis
$\{\Obar_a\}$, solve
\begin{equation}
  C(t_d)\, v \;=\; \lambda\, C(t_0)\, v ,
  \label{eq:gevp}
\end{equation}
take $v_0$ the eigenvector of the largest $\lambda$, and define the
\emph{projected operator} $\Obar_{v_0} = \sum_a (v_0)_a \Obar_a$. This is the
combination of basis members that maximises the decay ratio $C(t_d)/C(t_0)$; it
is what \texttt{gevp\_ground\_vector} returns and what every ``classical''
$\Azero$ in this report is measured on. Throughout, $t_0 = 1$ and $t_d = 2$.

\subsection{The whitening problem: flooring versus truncation}
\label{sec:whitening}

Equation~\eqref{eq:gevp} is solved by whitening: diagonalise $C(t_0) = Q s
Q^{\mathsf T}$, form $W = Q s^{-1/2}$, and take the top eigenvector of
$W^{\mathsf T} C(t_d) W$. When the basis is a long ladder of near-degenerate
operators --- seven smearing levels of the same plaquette are almost collinear,
and $21$ operators built as levels $\times$ shapes are worse --- $C(t_0)$
estimated on $400$ configurations is numerically singular and can be indefinite.
The $21$-operator arm reads $\operatorname{cond} C(t_0) = \infty$ on both
ensembles. There are two repairs:

\begin{description}
  \item[Flooring.] Clamp $s \to \max(s, \epsilon\, s_{\max})$ and keep all
        $n_{\text{ops}}$ directions. Every direction survives, including ones
        whose scale is set entirely by $\epsilon$.
  \item[Truncation.] Keep only the $k$ directions with $s > \epsilon\, s_{\max}$
        and \emph{drop} the rest, so the GEVP is solved in a $k$-dimensional
        subspace with $k \le n_{\text{ops}}$.
\end{description}

The difference is not cosmetic. A floored near-null direction is a direction
whose normalisation is arbitrary, and the whitened problem can and did select
one: an early run of this audit reported the strengthened arms as \emph{worse}
than the basis they contain --- $\Azero = 0.072$ for the $21$-operator arm
against $0.837$ for its own $4$-operator subset --- and duly announced that the
claim survived at $8.8\sigma$.\footnote{\texttt{notes/audit\_2026-09-06.md}
\S6. The failure mode is instructive precisely because it pointed the
flattering way.} Since \arm{published} is a strict \emph{subset} of the others,
their GEVP optimises over a larger span and cannot interpolate worse; the
violation of that superset property is what exposed the estimator. It is now a
hard gate in the script.

Because the choice of repair can move a number, it had to be made by a rule
fixed in advance rather than by which setting gave the nicer answer. That rule,
its failure and its amendment are Section~\ref{sec:selection}, and the evidence
is the four-setting table there.

\subsection{The decomposition: an inner product, not a fit}
\label{sec:innerproduct}

The third audit rests on an observation that is exact rather than approximate.
Operators acting on the vacuum are vectors in a Hilbert space, and the
equal-time connected correlator \emph{is} the inner product on that space:
\begin{equation}
  \langle O_a, O_b \rangle \;\equiv\; C_{ab}(0)
  \;=\; \big\langle\, d_a(t)\, d_b(t) \,\big\rangle ,
  \label{eq:metric}
\end{equation}
which is what \texttt{connected\_correlator\_matrix} already returns at
$\Delta = 0$. Consequently ``the part of the learned operator that the classical
basis cannot express'' is a \emph{literal orthogonal projection}:
\begin{equation}
  \OG \;=\; P + r,
  \qquad P = \operatorname{proj}_{\spn\{\Obar_a\}} \OG,
  \qquad r \perp \spn\{\Obar_a\} .
  \label{eq:decomp}
\end{equation}
Nothing here is fitted to come out any particular way. Three consequences make
the split worth measuring.

\begin{enumerate}
  \item \textbf{$P$ is the strongest possible classical opponent.} By
        construction it is the best approximation to $\OG$ \emph{inside the
        span} --- better than the GEVP vector, better than the most-smeared
        member, better than any classical operator chosen for any other reason.
        So $\dAzero(\text{GELT} - P) > 0$ is a \emph{lower bound} on the
        advantage over the entire span, and it is immune to the objection that
        answers Section~\ref{sec:fairfight} the hard way: one cannot claim the
        comparator was badly chosen when the comparator is the learned
        operator's own shadow.
  \item \textbf{The difference is better conditioned than the published one.}
        Equation~\eqref{eq:published} compares GELT to the GEVP vector, a
        \emph{different} operator correlated with GELT only through shared
        configurations. $P$ is GELT's own projection, so the two share every
        fluctuation living in the span and much of that variance cancels.
  \item \textbf{The same statistic gives the missing scale.} ``$12.9\%$ of the
        norm${}^2$ lies outside a four-dimensional span'' means nothing alone:
        any operator outside a $4$D subspace has a large orthogonal component.
        The scale is the basis's \emph{own} increments --- how much new content
        each rung of the smearing ladder adds to the span below it --- which
        costs no new data and says whether the ladder is still growing.
\end{enumerate}

\paragraph{The amplitude split, at one fixed mass.} $Z$ is linear in the
operator, so $Z_G = Z_P + Z_r$ holds \emph{exactly} --- but only if all three
are read against the same exponential. Letting each piece fit its own mass
destroys the additivity that makes the split interpretable, so amplitudes come
from a least-squares fit of $C(\Delta) \approx A\,\cosh_{\text{ref}}(\Delta)$ at
the learned operator's reference mass on that ensemble. The share $Z_r/Z_G$ is
then read off the \emph{cross}-correlator, $C_{Gr}(\Delta) \to Z_G Z_r
e^{-m\Delta}$ against $C_{GG}(\Delta) \to Z_G^2 e^{-m\Delta}$, which carries its
own sign and needs no square root. For a network whose own cosh fit does not
converge --- an untrained one --- the reference mass is passed in from the
trained net on the same ensemble (\texttt{-{}-m-ref}).

\subsection{Conventions, all pre-registered}
\label{sec:conventions}

Table~\ref{tab:conv} is the complete list. It was fixed in
\texttt{notes/fable5.1\_10-09\_audit.md} \S1.4 before the measurements, with the
single exception noted in its last row and explained in
Section~\ref{sec:selection}.

\begin{table}[htbp]
\centering\small
\caption{Fixed conventions. Every number in this report is measured under
these, on both ensembles, with no per-point tuning.}
\label{tab:conv}
\begin{tabular}{@{}lll@{}}
\toprule
item & value & where \\
\midrule
fit window            & $\Delta \in [2,7]$, cosh, profiled-$A$ grid & \texttt{fit\_glueball\_overlap.py} \\
jackknife             & delete-block, block $10$, whole fit per replica & same \\
GEVP                  & $t_0 = 1$, $t_d = 2$, $v_0$-projected operator & \texttt{gevp\_ground\_vector} \\
whitening             & \textbf{truncate}, $\epsilon = 10^{-4}$ relative & adopted, \S\ref{sec:selection} \\
test split            & contiguous last $20\%$ ($400$ configs) & \texttt{train\_glueball.py} \\
ensembles             & run5 (seed 0), ens1 (seed 1), inverse-variance & \S\ref{sec:intro} \\
lattice               & $SU(2)$, $12^3\times24$, $\beta=2.4$, $\xi=3.0$ & --- \\
width policy          & every new trained point at $d_{\text{model}} = 24$ & \S\ref{sec:width} \\
random trace          & init seeds $0,1,2$; error is the larger of seed & \S\ref{sec:randomtrace} \\
                      & spread and mean jackknife error & \\
\bottomrule
\end{tabular}
\end{table}

The two ensembles are sampled from scratch with different sampler seeds and
combined inverse-variance. They are not two halves of one ensemble: run5 and
ens1 have different classical anchors ($\meff(\Delta=1) \approx 0.357$ and
$0.390$) and different Rayleigh-loss floors, so a result replicating across them
has replicated across the ensemble noise, the dominant systematic here.

% ==============================================================================
\section{The fair fight: was the comparator a straw man?}
\label{sec:fairfight}
% ==============================================================================

\subsection{The arms}
\label{sec:arms}

Six classical bases are built on the same $400$ test configurations, projected
by the same \texttt{gevp\_ground\_vector}, fitted on the same window with the
same fixed $\sigma_\Delta$ and jackknifed on the same blocks --- the protocol is
\emph{imported} from \texttt{fit\_glueball\_overlap.py}, so ``same conventions
as the published table'' is a fact about the call graph. Each basis is a set of
spatial Wilson loops of shapes $R\times T$ built from links smeared with $n$
cumulative APE steps:

\begin{center}\small
\setlength{\tabcolsep}{4pt}
\begin{tabular}{@{}llcl@{}}
\toprule
arm & content & $n_{\text{ops}}$ & role \\
\midrule
\arm{thin}      & $n \in \{0\}$, shape $1\times1$ & $1$ &
   the unsmeared plaquette; the curve's left end \\
\arm{published} & $n \in \{0,2,4,6\}$, shape $1\times1$ & $4$ &
   the arm of Eq.~\eqref{eq:published} \\
\arm{deep}      & $n \in \{0,2,4,6,8,12,16\}$, shape $1\times1$ & $7$ &
   more radius; input-matched to the $7$-level net \\
\arm{shapes}    & $n \in \{0\}$, five $R\times T$ shapes & $5$ &
   extent without smearing \\
\arm{shapes\_sm}& $n \in \{0,2,4,6\}$, five shapes & $20$ &
   levels $\times$ shapes, shallow ladder \\
\arm{full}      & $n \in \{0,2,4,6,8,12,16\}$, three shapes & $21$ &
   the strongest basis available \\
\bottomrule
\end{tabular}
\end{center}

The loop shapes are cubic-symmetrised, so each ``shape'' is already a sum over
the orientations the $0^{++}$ channel requires. \arm{deep} is the
\emph{input-matched} comparator: its smearing levels are exactly the channels
the seven-level network was fed, so a win there cannot be explained by the
network having seen more smearing than the classical side. Two gates run before
any verdict: the \emph{slice gate} (the cached basis is measured on the same
configurations as the dump --- a failure means wrong configurations, not
physics) and the \emph{superset gate} (no superset arm reads below
\arm{published} by more than $2\sigma$, which by Section~\ref{sec:whitening}
cannot happen for a correct estimator, so a violation is a bug report and not a
result).

\subsection{The estimator table, and a pre-registered rule with no answer}
\label{sec:selection}

The rule fixed in advance was: run every arm on both seven-level dumps under
(a) floor at $\epsilon=10^{-4}$, (b) truncate at $\epsilon=10^{-4}$, (c)~(b)
plus collinearity pruning at $\rho = 0.999$, (d)~(b) plus pruning at
$\rho = 0.99$; then quote the \emph{least invasive setting under which every arm
passes both gates on both ensembles and \arm{published} is unchanged to three
decimals}. Sixteen offline runs, $45$ seconds in total on a laptop, produced
Table~\ref{tab:est}.

\begin{table}[htbp]
\centering\small
\caption{Classical $\Azero$ per arm under the four estimator settings, both
ensembles. ``kept'' is the number of $C(t_0)$ directions above
$10^{-4} s_{\max}$ under (b); ``prune'' the number of operators left by (c)/(d)
before truncation. Setting (c) reads (b) to $0.001$ everywhere and is omitted
from the $\Azero$ columns.}
\label{tab:est}
\begin{tabular}{@{}lcccll@{}}
\toprule
arm & $n_{\text{ops}}$ & kept (b) & prune (c/d)
    & $\Azero$ run5: (a) $\to$ (b) $\to$ (d)
    & $\Azero$ ens1: (a) $\to$ (b) $\to$ (d) \\
\midrule
\arm{thin}       & 1  & 1  & 1 / 1   & $0.390(60)$ & $0.484(128)$ \\
\arm{published}  & 4  & 4  & 4 / 4   & $0.837(56)$, unchanged & $0.925(71)$, unchanged \\
\arm{deep}       & 7  & 5  & 7 / 7   & $0.903 \to 0.905 \to 0.905$ & $0.939 \to 0.939 \to 0.939$ \\
\arm{shapes}     & 5  & 5  & 5 / 5   & $0.500(110)$ \quad \textsc{fell back} & $0.612(110)$ \quad \textsc{fell back} \\
\arm{shapes\_sm} & 20 & 11 & 20 / 16 & $0.934 \to 0.913 \to 0.910$ & $1.003 \to 1.004 \to 1.004$ \ \textsc{fb} \\
\arm{full}       & 21 & 9  & 20 / 14 & $0.942 \to 0.938 \to 0.936$ & $0.960 \to 0.960 \to 0.978$ \\
\bottomrule
\end{tabular}
\end{table}

\begin{table}[htbp]
\centering\small
\caption{Combined $\dAzero = \text{GELT} - \text{arm}$ (run5 $\oplus$ ens1,
inverse variance) under the four settings, for both trained networks. In
brackets, the run5\,--\,ens1 tension. ``same'' means identical to the quoted
digits.}
\label{tab:dest}
\begin{tabular}{@{}lllll@{}}
\toprule
arm & (a) floor & (b) truncate & (c) $+$ prune $0.999$ & (d) $+$ prune $0.99$ \\
\midrule
\multicolumn{5}{@{}l}{\emph{seven-level GELT}}\\
\arm{thin}        & $+0.569 \pm 0.069$ & same & same & same \\
\arm{published}   & $+0.131 \pm 0.029$ $(0.4\sigma)$ & same & same & same \\
\arm{deep}        & $+0.097 \pm 0.021$ $(1.7\sigma)$ & $+0.095 \pm 0.021$ $(1.8\sigma)$ & $+0.095 \pm 0.021$ & $+0.095 \pm 0.021$ \\
\arm{shapes}      & $+0.454 \pm 0.084$ & same & same & same \\
\arm{shapes\_sm}  & $+0.058 \pm 0.016$ $(0.7\sigma)$ & $+0.059 \pm 0.029$ $(0.2\sigma)$ & $+0.059 \pm 0.029$ & $+0.060 \pm 0.026$ \\
\arm{full}        & $+0.032 \pm 0.026$ $(1.5\sigma)$ & $\mathbf{+0.038 \pm 0.027}$ $(1.4\sigma)$ & $+0.037 \pm 0.026$ & $+0.042 \pm 0.022$ \\
\addlinespace
\multicolumn{5}{@{}l}{\emph{four-level GELT}}\\
\arm{thin}        & $+0.517 \pm 0.063$ & same & same & same \\
\arm{published}   & $+0.077 \pm 0.022$ $(0.5\sigma)$ & same & same & same \\
\arm{deep}        & $+0.041 \pm 0.025$ $(1.5\sigma)$ & $+0.039 \pm 0.025$ & $+0.039 \pm 0.025$ & $+0.039 \pm 0.025$ \\
\arm{shapes\_sm}  & $+0.002 \pm 0.024$ & $+0.006 \pm 0.035$ & $+0.006 \pm 0.035$ & $+0.008 \pm 0.032$ \\
\arm{full}        & $-0.022 \pm 0.026$ $(1.4\sigma)$ & $-0.013 \pm 0.030$ & $-0.013 \pm 0.030$ & $-0.011 \pm 0.023$ \\
\bottomrule
\end{tabular}
\end{table}

\paragraph{The gates pass.} \arm{published} reproduces $+0.066/+0.089$ against
the four-level networks --- i.e.\ Eq.~\eqref{eq:published} --- and its $\Azero$
is unchanged to the third decimal under all four settings. \arm{deep}
reproduces $+0.052/+0.127$ under (a) and reads $+0.050/+0.127$ under (b)--(d):
truncation drops $2$ of its $7$ directions and moves the answer by $0.06\sigma$.
The superset gate passes everywhere.

\paragraph{The rule has no answer as written.} Two arms, \arm{shapes} and
\arm{shapes\_sm}, fall back to their best single member on ens1 under
\emph{all four} settings, so no setting satisfies ``every arm passes on both
ensembles''. The interesting part is why, because the natural diagnosis ---
ill-conditioning --- is wrong: \arm{shapes} falls back at
$\operatorname{cond} C(t_0) = 87$, which is a perfectly healthy matrix.

Checked directly on the ens1 cache, the $v_0$-projected \arm{shapes} operator
\emph{beats} its best member on $C(2)/C(1)$, $0.640$ against $0.597$, and
\emph{loses} on $C(2)/C(0)$, $0.092$ against $0.206$. Both are correct and they
do not conflict: the GEVP at $(t_0,t_d) = (1,2)$ is by construction the
maximiser of $C(2)/C(1)$ and says nothing whatever about $C(2)/C(0)$ --- while
the variational gate, and $\Azero$ itself, are statements about
$C(\cdot)/C(0)$. The projection bought its improved $1\to2$ decay at the cost of
a large $\Delta=0$ piece ($C(1)/C(0) = 0.14$): GEVP-optimal and a poor $\Azero$
interpolator at once. \arm{shapes\_sm} on ens1 is the same story
($C(2)/C(1)$: $0.720$ vs $0.712$; $C(2)/C(0)$: $0.430$ vs $0.490$).

The consequence is structural, not numerical. The variational gate is a theorem
only at $t_0 = 0$, where the whitening metric and the normalisation of $\Azero$
are the same matrix; at $t_0 = 1$ they are different matrices, so no change to
the \emph{whitening} can rescue an arm that fails this way. A rule no setting in
the family can satisfy cannot select a setting.

\paragraph{The amendment, and why it moves nothing.} The rule was therefore
amended, and recorded as a deviation: it now ranges over the arms on the curve's
main chain (\arm{thin} $\subset$ \arm{published} $\subset$ \arm{deep} $\subset$
\arm{full}), since \arm{shapes} is not on the curve at all and
\arm{shapes\_sm} is the off-chain point, reported with its ens1 fallback flagged
rather than dropped. Among the settings that pass for those arms, the least
invasive one that removes \arm{full}'s null directions is (b): truncate at
$\epsilon = 10^{-4}$, no pruning --- the estimator for every number in this
report. The amendment changes no conclusion, and Table~\ref{tab:dest} is the
evidence rather than an assurance: every combined value there is stable across
(a)--(d) to $\le 0.01$, well inside its own error, and \arm{published} is
identical to three decimals under all four.

\subsection{The verdict on loop shapes}
\label{sec:shapesverdict}

With truncation, \arm{full} becomes readable: a nine-dimensional GEVP that
passes both gates on both ensembles, whose $\Azero$ barely moves from the floor
path ($0.942 \to 0.938$ on run5, $0.960 \to 0.960$ on ens1). This answers the
question the earlier audit left explicitly open.

\begin{equation}
  \dAzero(\arm{full}) = +0.038 \pm 0.027 \ (1.4\sigma)
  \qquad\text{against}\qquad
  \dAzero(\arm{deep}) = +0.095 \pm 0.021 \ (4.5\sigma).
  \label{eq:shapes}
\end{equation}

Loop shapes stacked on top of the deep ladder close about $60\%$ of the gap, and
what is left is a $1$--$2\sigma$ edge in the network's favour. That is precisely
the pre-registered expectation --- ``parity within a few hundredths or a
$1$--$2\sigma$ GELT edge'' --- which is worth saying because it is one of the
few predictions in this block that landed where it was aimed. The four-level
network, whose inputs contain neither the deep levels nor any shapes, is at
parity with \arm{full} ($-0.013 \pm 0.030$).

\begin{figure}[htbp]
\centering
\includegraphics[width=0.88\textwidth]{su2_fair_fight_best_glueball_gelt_sm0-2-4-6-8-12-16_trunc.png}
\caption{The fair fight on the run5 ensemble, seven-level network, adopted
(truncated) estimator. Top left: $\Azero$ of each classical arm against the
learned operator (horizontal line). Top right: the correlated difference
$\dAzero$ per arm on the same configurations, with the published $+0.078$
marked. Bottom left: $C(\Delta)/C(0)$ per arm --- \arm{thin} and \arm{shapes}
visibly lose their signal by $\Delta \approx 5$, which is why their $\Azero$ is
poor and why neither is a serious opponent. The bottom-right panel is the
$\mathbb{Z}_2$ smearing-freeze diagnostic, inactive here because this run reads
the cached basis rather than re-smearing an ensemble. Regenerate with
\texttt{SFF\_TRUNCATE=1 SFF\_BASES=1 python scripts/su2\_fair\_fight.py
dumps/best\_glueball\_gelt\_sm0-2-4-6-8-12-16\_test\_obars.pt}.}
\label{fig:fairfight}
\end{figure}

\paragraph{Two earlier statements that do not survive.} The first pass of this
audit, written when \arm{full} was unreadable, recorded two claims that the
truncated estimator refutes.\footnote{\texttt{notes/audit\_2026-09-06.md} \S6.5,
superseded on both points by \texttt{notes/fable5.1\_10-09\_audit.md} \S8.1.}

\begin{itemize}
  \item \emph{``\arm{full}'s $\dAzero$ is not reproducible between ensembles ---
        an $8\times$ swing.''} The two values, $+0.013 \pm 0.029$ and
        $+0.105 \pm 0.056$, differ by $1.4\sigma$. The $8\times$ is the ratio of
        two small central values, which is not a measure of tension.
  \item \emph{``\arm{full}'s eigenvector is fixed by the $\epsilon$ floor, not
        by signal.''} Dropping the floored directions entirely moves
        \arm{full}'s $\Azero$ by $0.004$. The floor was not steering the answer.
        What the earlier note got right is that $\operatorname{cond} = \infty$
        made that impossible to know without this check --- the objection was
        sound, the conclusion was not.
\end{itemize}

\paragraph{What this does and does not license.} \emph{Does}: the published
comparator was under-powered, and the honest headline against an input-matched
classical basis is $+0.095 \pm 0.021$ rather than $+0.078 \pm 0.022$; a
classical basis with real loop-shape variety closes most of the remaining gap,
and the network retains a $1$--$2\sigma$ edge over $21$ hand-built operators.
\emph{Does not}: treat the $21$-operator comparison as a measurement of equal
quality to the others. \arm{deep} and \arm{full} both sit at $\Azero \ge 0.90$
on both ensembles, inside the saturation band of Section~\ref{sec:a0}, so
Eq.~\eqref{eq:shapes} is read off a compressed observable and the $60\%$ is a
rough fraction, not a measured ratio with an error. Extending past that ceiling
needs a harder observable, not more configurations (Section~\ref{sec:next}).

% ==============================================================================
\section{The curve: architecture or inputs?}
\label{sec:curve}
% ==============================================================================

\subsection{Definition}
\label{sec:curvedef}

The fair fight strengthens the opponent while holding the network fixed. The
curve does something different: it varies the \emph{input content} $x$ handed to
\emph{both} methods and asks how the difference behaves. Write $\Azero(x)$ for
\begin{equation}
  x \;\in\; \{\ \text{thin},\ \text{4lv},\ \text{7lv},\ \text{4lv+5sh},\
  \text{7lv+3sh}\ \},
\end{equation}
ordered by inclusion along the main chain thin $\subset$ 4lv $\subset$ 7lv
$\subset$ 7lv+3sh, with 4lv+5sh an off-chain point (neither a subset nor a
superset of 7lv). Three traces are measured at each $x$:

\begin{center}\small
\setlength{\tabcolsep}{4pt}
\begin{tabular}{@{}lll@{}}
\toprule
trace & what it is & source \\
\midrule
\textbf{classical} GEVP$(x)$ & the $v_0$-projected GEVP at content $x$
  & the arm of that name \\
\textbf{trained} GELT$(x)$ & a GELT trained on channels with content $x$
  & \texttt{train\_glueball.py} \\
\textbf{random} GELT$(x)$ & the same architecture, untrained, same inputs
  & eval-only, three init seeds \\
\bottomrule
\end{tabular}
\end{center}

Every point is measured on the same $400$ held-out configurations, every
difference at fixed $x$ is a correlated blocked jackknife on shared
configurations, and both ensembles are combined inverse-variance. The third
trace is the load-bearing one for this section's title: if the architecture
alone, at random initialisation, already beat the classical GEVP of the same
inputs, ``learned'' would be the wrong word for what the trained network does.

\subsection{The measurements}
\label{sec:curvetable}

Twenty-six dumps, one fair-fight run each under the adopted estimator ($73$
seconds on a laptop for all of them), assembled by
\texttt{scripts/input\_architecture\_curve.py}, which fits nothing and
re-estimates nothing --- the numbers in Table~\ref{tab:curve} are the numbers
\texttt{su2\_fair\_fight.py} printed.

\begin{table}[htbp]
\centering\small
\caption{$\Azero$ against input content. ``classical'' is the GEVP over the same
input content; ``trained'' is at $d_{\text{model}} = 24$; ``random'' is the mean
over init seeds $0,1,2$ under the error rule of Table~\ref{tab:conv}. $\dAzero$
is the correlated trained\,--\,classical difference. The last column is the
rung-equivalence entry $\dAzero(\mathrm{GELT}(x_k) - \mathrm{GEVP}(x_{k+1}))$:
the learned operator at one input depth against the classical basis at the
next. Generated verbatim as
\texttt{results/fair\_fight/input\_architecture\_curve.tex}.}
\label{tab:curve}
\setlength{\tabcolsep}{4pt}
\begin{tabular}{lccccccc}
\toprule
$x$ & ens & $A_0$ classical & $A_0$ trained & $A_0$ random & $\Delta A_0$ & combined & rung $\Delta A_0$ \\
\midrule
thin & run5 & $0.390(60)$ & $0.651(77)$ & $0.249(62)$ & $+0.261 \pm 0.063$ & $+0.252 \pm 0.056$ & $-0.201 \pm 0.048$ \\
 & ens1 & $0.484(128)$ & $0.706(73)$ & $0.280(109)$ & $+0.221 \pm 0.117$ &  &  \\
\addlinespace
4lv & run5 & $0.837(56)$ & $0.911(51)$ & $0.497(175)$ & $+0.073 \pm 0.026$ & $+0.075 \pm 0.019$ & $+0.038 \pm 0.024$ \\
 & ens1 & $0.925(71)$ & $1.002(61)$ & $0.532(198)$ & $+0.078 \pm 0.029$ &  &  \\
\addlinespace
7lv & run5 & $0.905(62)$ & $0.955(57)$ & $0.730(343)$ & $+0.050 \pm 0.033$ & $+0.095 \pm 0.021$ &  \\
 & ens1 & $0.939(65)$ & $1.066(59)$ & $0.512(148)$ & $+0.127 \pm 0.027$ &  &  \\
\addlinespace
\bottomrule
\end{tabular}
\end{table}

Combined over ensembles, the classical trace reads $0.407 \pm 0.054$,
$0.870 \pm 0.044$, $0.921 \pm 0.045$; the trained trace $0.680 \pm 0.053$,
$0.949 \pm 0.039$, $1.008 \pm 0.041$; the random trace $0.257 \pm 0.054$,
$0.512 \pm 0.131$, $0.547 \pm 0.136$.

\begin{figure}[htbp]
\centering
\includegraphics[width=\textwidth]{input_architecture_curve.png}
\caption{The input--architecture curve. Left: $\Azero$ against input content for
the three traces, ensembles inverse-variance combined, with the two off-chain
classical points (\arm{shapes\_sm}, \arm{full}) as open triangles and the
$\Azero \ge 0.90$ saturation band shaded. The open star is the width control
(Section~\ref{sec:width}). Right: the correlated difference $\dAzero(x)$, and
the rung-equivalence difference $\mathrm{GELT}(x_k) - \mathrm{GEVP}(x_{k+1})$,
which is the quantity asking whether the architecture is worth a smearing rung.
Note the shape of the black curve: the advantage is large where the classical
basis is poor and shrinks as it is fed.}
\label{fig:curve}
\end{figure}

\subsection{The three readings, evaluated by their own tests}
\label{sec:readings}

A point is \emph{resolvable} when the combined classical $\Azero < 0.90$;
above that the observable has saturated and cannot discriminate. This excludes
7lv ($\Azero = 0.921$) and leaves thin ($0.407$) and 4lv ($0.870$). The three
readings and their pre-registered operational tests were fixed before any of
Table~\ref{tab:curve} existed, and the script evaluates them mechanically.

\begin{description}
  \item[R1 --- the architecture buys input depth.] \emph{Test:} $\dAzero(x) >
        2\sigma$ at every resolvable $x$, \textbf{and} rung-equivalence
        $\Azero^{\text{GELT}}(x_k) \ge \Azero^{\text{GEVP}}(x_{k+1})$ within
        errors for $k =$ thin$\to$4lv and 4lv$\to$7lv.
        \emph{Result: splits.} The first clause holds cleanly. The rung clause
        holds for 4lv$\to$7lv ($+0.038 \pm 0.024$: a four-level GELT matches the
        seven-level classical basis) and \textbf{fails} for thin$\to$4lv
        ($-0.201 \pm 0.048$: a thin-input GELT does not reach the four-level
        basis). R1 does not hold as stated.
  \item[R2 --- an efficient shortcut, not an extension.] \emph{Test:}
        $\dAzero(x)$ consistent with zero at some resolvable $x$ where the GEVP
        still has $1-\Azero$ more than $3\sigma$ from zero.
        \emph{Result: fails.} The traces never meet before saturation.
  \item[R3 --- an additive advantage.] \emph{Test:} $\dAzero(x)$ consistent with
        a constant $c > 0$ over the resolvable points, $\chi^2/\text{dof} < 2$
        and $c > 3\sigma$. \emph{Result: fails,} and informatively. The constant
        fit gives $c = +0.094 \pm 0.018$ with $\chi^2/\text{dof} = 9.1$: the
        advantage is not additive, it \emph{shrinks} as the inputs get richer,
        from $+0.25$ at thin links to $+0.08$ at four smearing levels.
\end{description}

That none of the three survives unmodified is the honest outcome and is more
informative than a clean pass would have been. What the data say is a fourth
statement the pre-registration did not name: the advantage is real at every
resolvable input depth, it is largest where the classical side is weakest, and
it decays toward the saturation ceiling rather than persisting as a fixed
offset. The rung clause splits along the same axis --- worth a whole smearing
rung at four levels, not worth one at thin links --- which is what one expects
if the network's contribution and the classical ladder's contribution overlap
increasingly as the ladder grows. Section~\ref{sec:decomp} measures that overlap
directly and finds exactly that: $21.1\% \to 12.8\% \to 2.5\%$ of the learned
operator lies outside the $4$-, $7$- and $21$-operator spans.

\subsection{The width control}
\label{sec:width}

The seven-level network is wider than the four-level one ($d_{\text{model}}$
$24$ against $16$), because it was trained with more input channels. A curve
drawn through both would then confound \emph{input content} with \emph{capacity}
--- the objection this section exists to answer, reappearing one level down. The
control is one extra training: the four-level input set at $d_{\text{model}} =
24$, on both ensembles.

\begin{center}\small
\begin{tabular}{@{}lcc@{}}
\toprule
4lv trained network & run5 & ens1 \\
\midrule
$d_{\text{model}} = 16$ (the published Run-5 net) & $0.903(47)$ & $1.013(62)$ \\
$d_{\text{model}} = 24$ (the curve's point)       & $0.911(51)$ & $1.002(61)$ \\
\bottomrule
\end{tabular}
\end{center}

Width $16 \to 24$ changes nothing on either ensemble. The curve is therefore
drawn at $d_{\text{model}} = 24$ throughout, so that moving along $x$ changes
the input content and nothing else, and the published $d_{\text{model}} = 16$
network becomes the control rather than a point. The $4\text{lv} \to 7\text{lv}$
gain is input content, not capacity.

\subsection{The random trace}
\label{sec:randomtrace}

The pre-registered expectation was $\Azero^{\text{random}}(x) \le
\Azero^{\text{GEVP}}(x)$ at every $x$, with the explicit note that a violation
anywhere would mean the architecture alone beats the optimal linear combination
of its own inputs, and that the ``learned'' attribution at that $x$ would then
rest on trained\,--\,random only.

It holds comfortably at all three points: $0.257$ against $0.407$ at thin links,
$0.512$ against $0.870$ at four levels, $0.547$ against $0.921$ at seven. An
untrained GELT is a legitimate gauge-invariant operator --- it has a mass and a
correlator --- but a distinctly poor one, well below the classical basis built
from the same configurations, so the advantage in Table~\ref{tab:curve} is
attributable to training rather than to the architecture alone. This is also the
first of the two controls \texttt{notes/operator\_decomposition.md} \S5 asked
for, delivered at three input depths rather than one.

\paragraph{What this does and does not license.} \emph{Does}: ``at every input
depth this observable can resolve, one learned operator carries more
ground-state weight than the optimal linear combination of the \emph{same}
inputs --- $+0.252 \pm 0.056$ on thin links, $+0.075 \pm 0.019$ at four smearing
levels; the advantage shrinks as the classical side is fed more, and at four
levels it is worth a whole extra smearing rung while at thin links it is not.''
\emph{Does not}: extend past 7lv. Above four levels the observable saturates and
cannot separate the methods further, and no amount of statistics changes that,
because the ceiling is $\Azero = 1$ and not a noise level. Nor does it license a
statement about the right-hand end of the chain: the shape-input network was
never trained (Section~\ref{sec:honesty}), so 7lv+3sh compares a GELT
\emph{without} loop shapes to a GEVP \emph{with} them, which is a fair reading
only as long as the figure says so --- it does, with that point drawn as an
off-chain classical marker and no trained counterpart.

% ==============================================================================
\section{The decomposition: what did the network find?}
\label{sec:decomp}
% ==============================================================================

Sections~\ref{sec:fairfight} and~\ref{sec:curve} establish that the learned
operator wins, and against what. Neither says what it \emph{is}. This section
does, using the exact orthogonal split of Section~\ref{sec:innerproduct}, and it
is where the report spends most of its numbers.

\subsection{The gate, and why it belongs in the report}
\label{sec:gate}

Before anything else is read, \texttt{operator\_decomposition.py} recomputes the
published comparison from \emph{its own} code path and prints it beside the
published values. It reproduces them to the last digit: $\Azero(\text{GELT}) =
0.9029 \pm 0.0468$ against the GEVP's $0.8374 \pm 0.0555$, $\dAzero = +0.0656
\pm 0.0306$ on run5 and $+0.0885 \pm 0.0302$ on ens1, $12.89(41)\%$ and
$11.69(30)\%$ of the norm${}^2$ outside the span, and a combined
$\dAzero(\text{GELT} - P) = +0.0761 \pm 0.0192$ ($4.0\sigma$). A second gate came
free when the span became selectable: with \texttt{-{}-basis=\dots:deep} the
GEVP arm reads $\Azero = 0.9053$ / $0.9390$ and $\dAzero(\text{GELT} -
\text{GEVP}) = +0.0496$ / $+0.1266$ --- Section~\ref{sec:fairfight}'s truncated
\arm{deep} row ($0.905/0.939$, $+0.050/+0.127$), reproduced from a script that
shares no estimator code with \texttt{su2\_fair\_fight.py} beyond the library fit
routine.

This belongs in the report and not only in the code, because it is the only
thing standing between a re-analysis and a silent discrepancy. The decomposition
runs on the same dumps as the fair fight but through different projection,
whitening and amplitude machinery; a disagreement of a few hundredths would be
invisible in a table of decomposition results alone, and every conclusion below
would inherit it. A gate that reproduces two independent published rows before
printing anything new converts that risk into a precondition, at the cost of one
extra column of an array the script already has.

\subsection{The main table: one operator against three nested spans}
\label{sec:spans}

The seven-level trained network is decomposed against three classical spans of
increasing size, combined over both ensembles. The columns are the three
statistics the split produces.

\begin{table}[htbp]
\centering\small
\caption{The \emph{same} trained seven-level operator against three nested
classical spans, both ensembles combined inverse-variance. \arm{full} is at the
script default \texttt{-{}-proj-eps} $=10^{-12}$; the last row repeats it at
$10^{-4}$, which retains $10$ of the $21$ Gram directions.}
\label{tab:spans}
\begin{tabular}{@{}lcccc@{}}
\toprule
span & $n_{\text{ops}}$ & norm${}^2$ outside & $Z_r/Z_G$
     & $\dAzero(\text{GELT} - P)$ \\
\midrule
\arm{published} (4lv)   & 4  & $0.2114 \pm 0.0040$ & $0.2293 \pm 0.0094$ & $+0.157 \pm 0.027$ \ $(5.8\sigma)$ \\
\arm{deep} (7lv)        & 7  & $0.1277 \pm 0.0025$ & $0.1459 \pm 0.0071$ & $+0.097 \pm 0.021$ \ $(4.5\sigma)$ \\
\arm{full} (7lv+3sh)    & 21 & $0.0251 \pm 0.0005$ & $0.0128 \pm 0.0038$ & $+0.021 \pm 0.009$ \ $(2.3\sigma)$ \\
\addlinespace
\arm{full}, $\epsilon = 10^{-4}$ & $21\!\to\!10$ & $0.0452 \pm 0.0010$ & $0.0374 \pm 0.0053$ & $+0.039 \pm 0.013$ \ $(2.9\sigma)$ \\
\bottomrule
\end{tabular}
\end{table}

\begin{table}[htbp]
\centering\small
\caption{Per-ensemble detail against \arm{deep}, the input-matched span. The
last two rows are the statements that matter: dropping $r$ costs the advantage,
and it does so at unchanged mass.}
\label{tab:deep}
\begin{tabular}{@{}lcc@{}}
\toprule
against $\spn\{\arm{deep}\}$ & run5 & ens1 \\
\midrule
norm${}^2$ fraction outside the span & $0.1273 \pm 0.0037$ & $0.1279 \pm 0.0033$ \\
$Z_r/Z_G$                            & $0.1287 \pm 0.0121$ & $0.1550 \pm 0.0088$ \\
$\Azero(P)$ --- best classical approximation of GELT & $0.8836 \pm 0.0578$ & $0.9444 \pm 0.0662$ \\
$\Azero(r)$ --- the remainder on its own & $0.4612 \pm 0.0697$ & $0.4165 \pm 0.0584$ \\
$\dAzero(\text{GELT} - P)$, correlated & $+0.0712 \pm 0.0304$ & $+0.1212 \pm 0.0301$ \\
$\Delta m(\text{GELT} - P)$            & $-0.0092 \pm 0.0133$ & $-0.0099 \pm 0.0142$ \\
\bottomrule
\end{tabular}
\end{table}

\paragraph{What each column means.}

\emph{The norm${}^2$ fraction outside the span} is $\|r\|^2/\|\OG\|^2$ in the
metric~\eqref{eq:metric} --- the squared sine of the angle between the learned
operator and the classical subspace. ``$12.8\%$ outside'' says that writing the
learned operator as a combination of the seven classical operators as well as
least squares allows leaves $12.8\%$ of its squared length unaccounted for. It
involves no fit and no mass, and --- crucially for
Section~\ref{sec:untrained} --- it says \emph{where the operator points}, not
whether it is any good.

\emph{$Z_r/Z_G$} is not a norm. It is the share of the \emph{ground-state
amplitude} carried by the residual, read off the cross-correlator at the fixed
reference mass (Section~\ref{sec:innerproduct}), so it carries a sign and it
asks a physical rather than a geometric question: of the operator's coupling to
the state we want, how much comes from the part the classical basis cannot
express? The two columns can and do disagree --- that disagreement is the entire
content of Section~\ref{sec:untrained}.

\emph{$\dAzero(\text{GELT}-P)$} is the sharp statement, and the one to quote. It
is not ``GELT beats a classical operator''; it is ``deleting from GELT
everything the classical span can express costs this much ground-state weight''.
Since $P$ is by construction the best classical approximation available, this is
a \emph{lower bound} on the advantage over the whole span. The accompanying
$\Delta m$, consistent with zero on both ensembles at $0.7\sigma$, says the two
are measuring the same state --- without which the $\Azero$ comparison would be
meaningless.

\subsection{Two readings that came out sideways}
\label{sec:sideways}

Two expectations were pre-registered before this measurement, and reporting how
they landed is more useful than reporting only the numbers.

\paragraph{``A smaller out-of-span fraction than the published $12.9\%/11.7\%$''
--- not as literally written.} Against \arm{deep}, the seven-level network reads
$12.7\%$ and $12.8\%$: the same number. At face value that is a strengthened
span absorbing nothing, which would be a serious problem for the whole study. It
is not, because the literal comparison mixes two different networks: the
published $12.9\%/11.7\%$ is the \emph{four-level} network against the
\emph{four-level} span. The apples-to-apples question --- does a stronger span
absorb the learned content? --- is one fixed operator against nested spans, and
there the answer is unambiguous:
\begin{equation}
  21.1\% \;\longrightarrow\; 12.8\% \;\longrightarrow\; 2.5\%
  \qquad (4 \to 7 \to 21 \text{ operators}).
  \label{eq:nested}
\end{equation}
Enlarging the span from four to seven levels absorbs $40\%$ of the learned
operator's out-of-span norm; adding loop shapes absorbs most of the rest. The
coincidence of $12.7\%$ with $12.9\%$ is exactly that --- a coincidence between
two different measurements. What it does reveal, read correctly, is that the
seven-level network is \emph{further} outside the four-level span ($21.1\%$)
than the four-level network is ($12.9\%$): richer inputs produce a more distant
operator. A study that quoted only the headline fraction would have reported the
same number twice and concluded nothing.

\paragraph{``$\dAzero(\text{GELT}-P)$ positive but below $+0.097$'' --- positive,
and \emph{at} the bound.} Measured: $+0.0965 \pm 0.0214$, against the fair
fight's $\dAzero(\text{GELT} - \arm{deep}\ \text{GEVP}) = +0.097 \pm 0.021$. The
two agree to the third decimal. The expectation was that they would not: $P$ is
a better classical operator than the GEVP vector, so deleting $r$ should cost
\emph{less} than losing to the GEVP.

That they coincide is not a coincidence, and the mechanism is visible in
Table~\ref{tab:deep}: $\Azero(P)$ lands on $\Azero(\text{GEVP})$ on both
ensembles --- $0.884$ against $0.905$ on run5, $0.944$ against $0.939$ on ens1.
The least-squares shadow of the learned operator and the variationally optimal
classical operator are the same quality of interpolator. That is a non-trivial
internal consistency check: the span really does have a ceiling near
$\Azero \approx 0.9$, and \emph{the ceiling is a property of the span rather
than an artifact of how the GEVP picks its vector}. The same check was made at
four levels in the original study ($0.830$ against $0.837$) and it now holds at
seven.

\subsection{What $r$ is made of}
\label{sec:shapespan}

The residual $r$ is orthogonal to $\spn\{\arm{deep}\}$ by construction, but it
need not be orthogonal to everything classical. Since
$\spn\{\arm{full}\} \supset \spn\{\arm{deep}\}$ and the extra directions are
exactly the loop shapes, projecting $r$ onto $\spn\{\arm{full}\}$ isolates how
much of the learned residual is rectangular Wilson loops:
\begin{equation}
  \frac{\|\operatorname{proj}_{\spn\{\arm{full}\}} r\|^2}{\|r\|^2}
  \;=\; 75.8(4)\%\ \text{(run5)},\quad 81.9(3)\%\ \text{(ens1)},
  \quad \text{combined } 79.6 \pm 0.3\%.
  \label{eq:shapefrac}
\end{equation}
This number depends on how much of \arm{full}'s near-degenerate span one is
willing to call real. At the script default the Gram of the $21$-operator span
is inverted with a relative cut of $10^{-12}$, keeping all $21$ directions and
giving Eq.~\eqref{eq:shapefrac}. Cutting at $10^{-4}$, which keeps $10$
directions, gives $55.5(7)\%$ and $60.7(8)\%$.\footnote{Both ends are quoted in
\texttt{notes/fable5.1\_10-09\_audit.md} \S8.3. That section describes the
$10^{-4}$ setting as truncating \arm{full} ``to the 10 directions \S8.1's
estimator keeps''; strictly, the $10$ directions are those of the \emph{$C(0)$
Gram} of the span under the same $\epsilon$, while \S8.1's estimator keeps $9$
directions of $C(t_0)$. The two are different matrices, so both counts are
correct, but they are not the same $10$ and $9$ directions.} The spread is the
honest error on this statement, and it is larger than the jackknife error: it is
a systematic of the projection, not of the sample.

The reading, at either end: \textbf{most of what the network found outside the
smearing ladder is rectangular loops it rediscovered} --- between a half and
four fifths --- and the remainder, about $2.5\%$ of the operator's total
norm${}^2$, is content that \emph{no} operator in the $21$-member classical
basis expresses. That last piece is not negligible: it still carries
$\dAzero = +0.021 \pm 0.009$ ($2.3\sigma$), consistent with the fair fight's
independent $\dAzero(\text{GELT} - \arm{full}) = +0.038 \pm 0.027$.

\subsection{The untrained control, and why it is the whole point}
\label{sec:untrained}

Everything above is compatible with a deflationary explanation: the L1-ball
parallel transport of the GELT block reaches offsets and path-averages that no
smeared plaquette contains, so \emph{any} network of this architecture ---
trained or not --- would sit outside the classical span. If that is the case,
the norm fraction is a fact about the architecture and carries no evidence of
learning at all. This was written down as the expectation, in those words,
before the control ran.

Six untrained seven-level networks (init seeds $0,1,2$ on both ensembles),
evaluated through the same pipeline against the same \arm{deep} span, with the
reference mass for the amplitude split taken from the trained network on that
ensemble ($0.3375$ on run5, $0.3689$ on ens1).

\begin{table}[htbp]
\centering\small
\caption{Trained against untrained, same architecture, same inputs, same span,
same $400$ configurations. Seed $0$'s cosh fit does not converge on either
ensemble ($m = 1.23 \pm 0.43$ and $0.75 \pm 0.30$ against the trained $0.34$ and
$0.37$), so it is excluded on a criterion visible without looking at $\dAzero$;
both combinations are reported.}
\label{tab:untrained}
\begin{tabular}{@{}lccc@{}}
\toprule
 & trained (7lv) & untrained (4 converged seeds) & untrained (all 6) \\
\midrule
norm${}^2$ outside span        & $0.1277 \pm 0.0025$ & $0.2371 \pm 0.0551$ & $0.4036 \pm 0.1220$ \\
$Z_r/Z_G$                      & $0.1459 \pm 0.0071$ & $0.0474 \pm 0.0220$ & $0.1894 \pm 0.1037$ \\
$r$ inside $\spn\{\arm{full}\}$& $0.7959 \pm 0.0027$ & $0.0028 \pm 0.0010$ & $0.0042 \pm 0.0010$ \\
$\dAzero(\text{net} - P)$      & $\mathbf{+0.0965 \pm 0.0214}$ & $\mathbf{-0.0911 \pm 0.0420}$ & $-0.0745 \pm 0.1271$ \\
\bottomrule
\end{tabular}
\end{table}

Per seed, $\dAzero(\text{net}-P) = -0.107(61)$ and $-0.064(35)$ on run5 and
$-0.197(60)$ and $-0.024(45)$ on ens1 for seeds $1$ and $2$; seed $0$ gives
$+0.70 \pm 0.91$ and $-0.13 \pm 0.31$, i.e.\ no information in either direction.
The per-seed out-of-span fractions are $16\%$, $31\%$ and $74\%$ (seeds $2$, $1$,
$0$), and the two ensembles agree to better than $1\%$ relative on each seed ---
as they must, since the init seed \emph{is} the network and the ensemble only
supplies configurations to evaluate it on. Three conclusions follow, the first
two pre-registered and the third not.

\begin{enumerate}
  \item \textbf{Being outside the span is architectural.} Untrained networks are
        $16$--$74\%$ outside, against the trained network's $12.8\%$. Training
        does not push the operator away from the classical span; it moves it
        \emph{closer} while making it better. The norm fraction therefore proves
        nothing about learning, exactly as the pre-registration said it would
        not, and any write-up that leads with ``$12.9\%$ outside the span'' as
        evidence of what the network learned is making an unsupported claim.
  \item \textbf{The discriminating statistic is $Z_r/Z_G$.} Per unit of norm,
        the trained residual is an order of magnitude more productive:
        \begin{equation}
          \frac{Z_r/Z_G}{\sqrt{\|r\|^2/\|\OG\|^2}}
          \;=\; \frac{0.146}{\sqrt{0.128}} = 0.41
          \quad\text{trained},
          \qquad
          \frac{0.047}{\sqrt{0.237}} = 0.10
          \quad\text{untrained}.
        \end{equation}
        An untrained network has a large out-of-span component that couples to
        almost nothing; training builds a smaller one that couples strongly.
  \item \textbf{The sign of $\dAzero(\text{net}-P)$ flips.} For the trained
        network, deleting the out-of-span part costs the whole advantage,
        $+0.097 \pm 0.021$. For an untrained one, deleting it \emph{helps}: its
        own projection onto the classical span is the better operator, by
        $0.091 \pm 0.042$. This is the cleanest separation in the report,
        because it is a change of sign rather than of magnitude, and no
        normalisation convention can produce it.
\end{enumerate}

\begin{figure}[htbp]
\centering
\includegraphics[width=0.94\textwidth]{operator_decomposition_deep_7lv.png}\\[0.8em]
\includegraphics[width=0.94\textwidth]{operator_decomposition_deep_7lv_rnd_fitok.png}
\caption{The decomposition against $\spn\{\arm{deep}\}$. \textbf{Top:} the
trained seven-level network. Left, the ladder increments of
Section~\ref{sec:ladder} with the learned operator as a star; right,
$\rho(\Delta) = [C(\Delta)/C(0)]/\cosh_{\text{ref}}(\Delta)$ for the three
pieces --- flat means pure, and the excess at small $\Delta$ is excited-state
contamination. $\OG$ (blue) sits above its own classical shadow $P$ (orange)
across the whole fit window; $r$ (red) is flat from $\Delta \approx 3$, i.e.\ it
couples to the same ground state rather than dying with the UV.
\textbf{Bottom:} two untrained networks (seeds $1$ and $2$; the left panel's two
stars are their out-of-span fractions, $31\%$ and $16\%$). The ordering in the
right panel is \emph{inverted} --- the projection $P$ lies above the network's
own operator --- which is the sign flip of Section~\ref{sec:untrained} drawn
rather than tabulated. (The bottom-left legend labels both traces ``Run 5''; the
classical ladder is in fact ensemble-specific and the two curves coincide to
line width, which is the cosmetic reason they are indistinguishable.)}
\label{fig:decomp}
\end{figure}

The third conclusion was not anticipated anywhere and is arguably the strongest
of the three: \textbf{the residual is loop-shaped only when the network is
trained.} $80\%$ of the trained $r$ is reachable inside $\spn\{\arm{full}\}$;
$0.3\%$ of the untrained $r$ is. An untrained equivariant network sits outside
the classical span in a direction that is outside \emph{everything} classical
--- some accident of the random transport average, orthogonal to the space
rectangular Wilson loops span. Training rotates that content onto the part of
operator space the loops do span, and then keeps a small piece that is not. That
is a much more specific description of what training does than ``it improves
$\Azero$''.

\subsection{The mechanism: constructive interference, not a better operator}
\label{sec:mechanism}

The naive reading of Section~\ref{sec:spans} --- ``the network found a better
operator hiding outside the basis'' --- is wrong, and the data say so directly.
On its own, $r$ is a \emph{poor} operator: $\Azero(r) = 0.46$ and $0.42$ on the
two ensembles, against the full learned operator's $0.95$ and $1.07$
(Table~\ref{tab:deep}); the published four-level study read $0.43$ and $0.37$.
It does not win by being good. It wins by adding \emph{coherently}.

The arithmetic is the reason the amplitude split is taken at one fixed mass.
With $Z_G = Z_P + Z_r$ exactly, removing $r$ removes a fraction
\begin{equation}
  1 - \left(1 - \frac{Z_r}{Z_G}\right)^{2}
  \label{eq:interference}
\end{equation}
of the ground-state \emph{intensity}. For the published four-level decomposition
$Z_r/Z_G = 0.131$, so Eq.~\eqref{eq:interference} gives $24.5\%$ of the intensity
removed by deleting $12.9\%$ of the norm${}^2$; against \arm{deep},
$Z_r/Z_G = 0.146$ gives $27\%$ from $12.8\%$ of the norm${}^2$. In both cases
the amplitude share is about twice what the norm share would suggest, the
ground-state fraction $\Azero = A/C(0)$ falls accordingly, and that is the whole
effect. (The numbers are of the right size to account for the observed drops ---
$0.903 \to 0.830$ at four levels, $0.955 \to 0.884$ at seven --- but are not
expected to match to three digits, because each arm's $\Azero$ is quoted from
its own cosh fit rather than from the fixed-mass amplitudes.)

Stated carefully for the thesis: what the network found is a \emph{direction}
that is mediocre alone and valuable in superposition --- precisely what a
variational method exists to exploit, and precisely what a basis of four, seven
or twenty-one hand-built operators cannot reach unless somebody thinks of the
right operator to add.

\subsection{The scale: the classical ladder's own increments}
\label{sec:ladder}

The remaining question from Section~\ref{sec:innerproduct} is what $13\%$ means.
The calibration is the same statistic applied to the classical operators
themselves: how much new content does each rung of the smearing ladder add to
the span of the rungs below it?

\begin{center}\small
\begin{tabular}{@{}lcc@{}}
\toprule
addition & run5 & ens1 \\
\midrule
APE$\times2$ outside $\spn\{\times0\}$ & $0.6891 \pm 0.0098$ & $0.7049 \pm 0.0084$ \\
APE$\times4$ outside $\spn\{\times0,2\}$ & $0.1971 \pm 0.0047$ & $0.1957 \pm 0.0043$ \\
APE$\times6$ outside $\spn\{\times0,2,4\}$ & $0.0486 \pm 0.0014$ & $0.0500 \pm 0.0012$ \\
APE$\times8$ outside $\spn\{\times0,2,4,6\}$ & $0.0120 \pm 0.0004$ & $0.0124 \pm 0.0003$ \\
APE$\times12$ outside $\spn\{\times0,\dots,8\}$ & $0.0207 \pm 0.0006$ & $0.0208 \pm 0.0006$ \\
APE$\times16$ outside $\spn\{\times0,\dots,12\}$ & $0.0098 \pm 0.0003$ & $0.0098 \pm 0.0003$ \\
\midrule
\textbf{GELT outside $\spn\{\text{the whole 7-level ladder}\}$}
  & $\mathbf{0.1273 \pm 0.0037}$ & $\mathbf{0.1279 \pm 0.0033}$ \\
\bottomrule
\end{tabular}
\end{center}

Each number replicates across two independently sampled ensembles to the third
decimal --- these are properties of the operator basis, not of the sample. The
first three rungs fall by a factor of $3.5$--$4$ each, and the learned
operator's $12.8\%$ exceeds every increment past the second: a seventh smearing
level buys $1\%$ of new content, the network buys $13\%$.

\paragraph{One extrapolation is withdrawn.} The original study fitted the decay
of the first three rungs and extrapolated that the \emph{entire} remaining tail
of the ladder --- rungs $8$, $12$, $16$ and beyond --- was worth $\approx 2\%$ of
new content, concluding that ``the classical smearing ladder is a geometric
series that has essentially converged''. The three rungs have since been
measured directly and they are $1.20\%$, $2.07\%$, $0.98\%$: a tail worth
$\approx 4\%$, not $\approx 2\%$, and \textbf{not monotone in the rung index}.
The non-monotonicity is not noise --- the errors are at the fourth decimal and
both ensembles show the same rise at $\times12$ --- it is that the ladder's own
steps are $\Delta2, \Delta4, \Delta4$, so the rungs are not equally spaced and
the series was never going to be geometric in the rung index. The extrapolation
is withdrawn; the qualitative reading (the ladder has slowed, the network is not
on that curve) survives, but ``the ladder has converged'' is too strong. It has
slowed, and its tail still matters at the percent level that this whole study
works in.

The same arithmetic underwrote a prediction the decomposition made about the
fair fight before that fight was readable: \arm{deep} could not close a $13\%$
gap, because more smearing is the direction the ladder has already converged
along, and only real loop geometry could move $\Azero$ materially. Both halves
held --- \arm{deep} leaves $+0.095 \pm 0.021$, and \arm{full} is where the
movement is --- but the numbers behind the prediction (rungs $8$, $12$, $16$ at
$1.3\%$, $0.35\%$, $0.1\%$) were too optimistic by the factors above. The
conclusion was right for a partly wrong reason, and both halves of that should
be said.

% ==============================================================================
\section{Honesty section: what these three studies do not settle}
\label{sec:honesty}
% ==============================================================================

\paragraph{The pre-registration, and what it bought.} The observable, the
estimator family, the fit conventions, the measurement grid and the three
readings with their tests were committed to git before any of the numbers in
Sections~\ref{sec:fairfight}--\ref{sec:decomp} existed. What that bought is
narrow: it removes the objection ``the opponent was chosen after the fact''. It
does \emph{not} make the results stronger than their error bars, and it did not
stop two of the three readings from failing their own tests. Two limits apply.
The estimator rule had to be amended (Section~\ref{sec:selection}); the
amendment is recorded with its reason and the four-setting table shows it moves
nothing, but an amended rule is weaker than an unamended one. And the
decomposition study was originally written \emph{after} its own first
measurement --- its window and metric were inherited rather than chosen, both
are scanned rather than reported at one setting, and the two ensembles were
analysed with identical code and agree, none of which is as good as
pre-registration and should not be described as if it were. The two
\emph{controls} in that section \emph{were} pre-registered, before they ran.

\paragraph{The seven-level network is a reactive point.} It was trained
\emph{after} the \arm{deep} arm was seen to win on the run5 ensemble. The
classical arms themselves were named before the fair fight ran and that pedigree
is real, but the network at the right-hand end of the main chain was a response
to a number. The thesis must say so where it quotes the seven-level numbers.
Everything else on the curve --- the thin retrain, the width control, the random
trace, the estimator and the readings --- was fixed before its numbers existed.

\paragraph{Two expectations came out sideways}, and both are in
Section~\ref{sec:sideways} rather than smoothed over: the out-of-span fraction
against the strengthened span did not shrink as literally predicted (it is the
same number, for a reason that turns out to be about comparing two different
networks), and $\dAzero(\text{GELT}-P)$ landed \emph{at} the bound rather than
below it (for a reason that is an internal consistency check in disguise). A
write-up that reported only the numbers would have lost both.

\paragraph{Saturation, a compressed observable, and $\Azero > 1$.}
$\Azero \le 1$ and every strong arm reads $0.90$--$1.07$. The 7lv point is
excluded from the curve's tests on that basis, and --- less comfortably --- the
whole of Section~\ref{sec:shapesverdict} is read off arms inside the saturation
band, so ``loop shapes close about $60\%$ of the gap'' is a ratio of two small
differences near a ceiling: the right qualitative reading, and not a number to
quote with an error bar. Relatedly, the learned operator reads
$\Azero = 1.013 \pm 0.062$ at four levels and $1.066 \pm 0.059$ at seven on ens1.
That is a fit artifact of a nearly pure operator, not a probability violation
(Section~\ref{sec:a0}); it is inherited from the published table, not made worse
here, and reported unclipped so that $1 - \Azero$ keeps its sign in every
difference. The thesis should say so once rather than leave a reader to notice
it.

\paragraph{The contact-term test does not survive a longer ladder.} The original
study repeated the whole decomposition under the metrics $C_{ab}(\tau)$ for
$\tau = 1, 2$, which weight the light states more heavily: a contact term --- UV
content living in $C(0)$ and gone by $\Delta = 2$ --- must \emph{shrink} as
$\tau$ grows. On the four-operator span it \emph{grew} ($0.129 \to 0.156 \to
0.197$ on run5), which is the opposite of a contact term and is corroborated by
$\rho(\Delta)$ for $r$ being flat from $\Delta \approx 3$
(Figure~\ref{fig:decomp}, top right). That test does not transfer. On
$\spn\{\arm{deep}\}$ the $C(2)$ Gram has $\lambda_{\min} = -2.9\times10^{-5}$
and on $\spn\{\arm{full}\}$ it is $-0.73$: $C(\tau)$ is then not an inner
product on those spans, the ``projection'' is undefined, and the jackknife error
duly explodes (to $\pm 383$ on run5). The script now checks the spectrum and
prints \texttt{NOT USABLE} beside the row rather than a number. \textbf{The
contact-term conclusion stands only where it was measured} --- the four-operator
span, where $C(1)$ and $C(2)$ are positive definite --- and is not restated for
the longer ladders.

\paragraph{The shape-span fraction depends on the projection cut.} $80\%$ at
\texttt{-{}-proj-eps} $=10^{-12}$ and $58\%$ at $10^{-4}$
(Section~\ref{sec:shapespan}). The $21$-operator span's $C(0)$ Gram has
condition number $6.9\times10^{8}$, so a cut is unavoidable; which value is
``right'' is a judgement about how much of a near-degenerate span is real. Both
ends are quoted everywhere the number appears. Nothing in the \arm{deep}
decomposition moves with this cut ($0.1273 \to 0.1276$ and $+0.0712 \to +0.0717$
from $10^{-12}$ to $10^{-4}$); only the shape fraction does.

\paragraph{One untrained seed is excluded.} Seed $0$'s cosh fit does not
converge on either ensemble ($m = 1.23 \pm 0.43$ and $0.75 \pm 0.30$ against the
trained $0.34$ and $0.37$; $\Azero = 1.02 \pm 0.89$ and $0.37 \pm 0.25$). It is
excluded on a criterion --- ``the fit did not converge'' --- visible without
looking at $\dAzero$, which is the only kind of exclusion that is safe here, and
\emph{both} combinations appear in Table~\ref{tab:untrained}. Including it does
not change the sign of $\dAzero(\text{net}-P)$; it inflates the error threefold,
which is what an unconverged fit should do.

\paragraph{The right-hand end of the curve is not matched, and one point is off
width.} The shape-input network was never trained, so the 7lv+3sh point compares
a GELT \emph{without} loop shapes to a GEVP \emph{with} them; the figure draws
it as an off-chain classical marker with no trained counterpart, and the fair
fight's $\dAzero(\arm{full}) = +0.038 \pm 0.027$ is the number that reading
rests on (closing it is Section~\ref{sec:next}, item~5). Separately, the
untrained trace at 4lv is at $d_{\text{model}} = 16$ rather than the curve's
$24$; six eval-only runs would fix it, and no reading depends on it, since that
trace is far below the classical arm at every $x$ and the width control
(Section~\ref{sec:width}) finds the two widths indistinguishable.

\paragraph{One theory, one lattice family, one architecture.} $SU(2)$,
$12^3\times24$, $\beta = 2.4$, $\xi = 3$, receptive radius $R = 2$, the Run-5
hyperparameters, the $0^{++}$ channel. There is no continuum limit here and none
is claimed. Every statement in this report is about \emph{this} learned operator
against \emph{these} classical bases on \emph{these} two ensembles.

% ==============================================================================
\section{What to do next, and why}
\label{sec:next}
% ==============================================================================

The items below are ranked by value per unit of cost, and each is stated as what
it measures, why the results above make it the right next thing, what it costs,
and what it would license.

\paragraph{1. The excited-state gate (laptop, minutes, offline).} Every
statement in this report is compressed against the $\Azero \le 1$ ceiling
(Section~\ref{sec:honesty}), and the only way past it is a harder observable ---
the obvious candidate being the first excited state, with the network emitting a
\emph{vector} of operators and a GEVP run on the learned pair. Before any
training code is written, one question must be answered: does a classical
comparator for such a basis exist on these ensembles at all? That is a pass/fail
test costing minutes, offline, from the tracked
\texttt{dumps/su2\_fair\_fight\_obars\_run5.pt}: solve the GEVP over
$\spn\{\arm{deep}\}$, read the \emph{second} eigenvalue's effective mass at
$\Delta = 1$ and $2$ under a delete-block jackknife, and pass only if the two
agree within $1\sigma$ with the $\Delta=2$ error below $25\%$ on at least one
ensemble. If it fails the option closes for free; if it passes, the design is
already fixed ($\texttt{mlp\_out} = 2$, loss $-\Tr[(C(0)+\varepsilon)^{-1}C(1)]$
over the learned $2\times2$ block, which by Ky Fan is maximised at
$e^{-m_1}+e^{-m_2}$ and penalises collapse onto two copies because $C(0)$ goes
singular). The cheapest item in the plan, and the only one gating a training
run.

\paragraph{2. Write the thesis text (no compute).} Every number the
spectroscopy chapter was waiting on now exists. Concretely: the single
``$+0.078 \pm 0.022$'' sentence becomes the curve statement of
Section~\ref{sec:readings}, with the four-level number kept as the first result
and the input-matched $+0.095 \pm 0.021$ as the headline; the overlap table
becomes the per-$x$ table; and the fair fight and the decomposition, which the
thesis does not currently mention at all, get a paragraph each. The
decomposition paragraph can now say \emph{learned} --- but on $Z_r/Z_G$ and the
sign flip of Section~\ref{sec:untrained}, never on the norm fraction, which the
untrained control shows to be architectural. This is the highest-value item that
costs nothing but writing, and doing it before the next measurement is what
keeps the measurements honest.

\paragraph{3. Fix the $\mathbb{Z}_2$ smearing ($\approx12$\,h GPU).} The one
live defect in the repository. At the production $\alpha = \tfrac12$ the APE
update for two spatial staples is $V = (1-\alpha)U + (\alpha/2)(s_1+s_2)$, which
is exactly $0$ whenever the two staples disagree with the link --- and
$\mathbb{Z}_2$'s projection sends $0 \to +1$, a value that does not transform
under a gauge transformation. Measured on production configurations the ladder
freezes after one step, $99.9\%$ of the changed links are gauge-dependent
tie-breaks, and the timeslice operator moves by $1.82\sigma$ under a gauge
transformation. The consequence for the tables in the attention report is that
the $\mathbb{Z}_2$ classical basis is \emph{one} operator and not four (the
dumps already record \texttt{gevp\_fell\_back=True, n\_ops=1}), and the
$\mathbb{Z}_2$ networks were trained on four channels of which three are
byte-identical. The fix is unprojected fat links --- linear in the links, hence
exactly covariant --- with the cheap first step being to evaluate the existing
checkpoints on covariant inputs before retraining anything. $SU(2)$ is verified
clean at $1.4\times10^{-15}$, so nothing in \emph{this} report is affected; the
next item, however, depends on it.

\paragraph{4. Dump per-configuration $\Obar$ from the $\mathbb{Z}_2$ attention
script (small code change).} The tooling built for Section~\ref{sec:decomp} ---
selectable spans, the shape-span projection, the fixed reference mass, the
untrained control --- is generic: it takes per-config $\Obar$ arrays and a
classical basis, and the attention field is exactly such an operator. Dumping
those arrays transports the whole analysis to the attention field with no GPU
re-run, turning ``the attention field has a mass and a good overlap'' into
``here is the part of it the classical basis cannot express''. The catch is item
3: until the $\mathbb{Z}_2$ smearing is fixed that classical span is
\emph{one-dimensional}, and a decomposition against a one-dimensional span is
not worth running. Sequence the two together, in this order.

\paragraph{5. Train a shape-input network (2 trainings).} Optional before
Section~\ref{sec:shapespan}, this is now the direct test of its main finding. If
$80\%$ of what the network found outside the smearing ladder is rectangular
loops it rediscovered, handing it loop shapes as input channels should let it
spend its capacity elsewhere --- and the prediction is sharp: the out-of-span
fraction against $\spn\{\arm{full}\}$ need not simply drop, since the network is
now \emph{given} the directions it previously had to build. It also closes the
unmatched right-hand end of the curve (Section~\ref{sec:honesty}). Cost:
cubic-symmetrised matrix-valued $R\times T$ loops as extra adjoint channels,
about a day of code, then two trainings of roughly eight hours each.

\paragraph{6. The $SU(2)$ $\beta$-scan of the attention field ($\approx6$\,h
GPU, zero code).} The attention report's $SU(2)$ row is a single coupling,
because $\beta = 2.4$ is the only one with a trained network, and the script
refuses to print aggregates below three couplings rather than emit a meaningless
correlation. Sampling four more ensembles and evaluating the existing network
off-diagonal turns that row into a scan, making the $\xi_A$ claim in $SU(2)$ the
same kind of statement it already is in $\mathbb{Z}_2$. Pure GPU filler with no
code to write --- and it must not share the GPU with a training, which is the
only reason it is not higher.

\paragraph{7. The matched-parameter L-CNN ($2$ trainings, plus a model switch).}
The one baseline the thesis names and has never run. The argument for attention
currently rests on beating a \emph{non}-equivariant CNN, which is a weak
opponent for the same reason \arm{published} was: it is not the thing a
practitioner would build. \texttt{gelt/lcnn.py} already mirrors GELT's I/O, so
the change is a model switch and axis-aligned transports in place of the L1-ball
average. The right presentation is not a standalone number but a \emph{fourth
trace} on Figure~\ref{fig:curve}: the same $\Azero(x)$ for a matched-parameter
L-CNN. The pre-registered expectation is parity with GELT --- the smeared inputs
hand both networks the staple content, and at matched depth the loop degree is
the same --- and parity is the honest form of ``attention costs nothing, and the
attention field is a measurement a convolution cannot produce''.

% ==============================================================================
\section{Conclusions}
\label{sec:conclusions}
% ==============================================================================

One number, $\dAzero = +0.078 \pm 0.022$, has been replaced by three answers to
three different questions, all measured offline on the same held-out
configurations under conventions fixed in advance.

\begin{itemize}
  \item \textbf{The comparator was under-powered, and strengthening it does not
        overturn the claim.} Against the input-matched seven-level ladder the
        learned operator leads by $+0.095 \pm 0.021$; against the strongest
        $21$-operator basis the repository builds, by $+0.038 \pm 0.027$. Loop
        shapes --- not more smearing --- close most of the gap, as the
        decomposition predicted before that arm was readable. Reading it at all
        required truncating the whitening rather than flooring it, and the
        pre-registered selection rule had to be amended, because two arms fail a
        gate the GEVP at $t_0 = 1$ does not guarantee: the projection maximises
        $C(2)/C(1)$ while the gate tests $C(2)/C(0)$.
  \item \textbf{The advantage is training, not architecture, and not merely
        richer inputs.} At every input depth the observable can resolve, one
        learned operator carries more ground-state weight than the optimal
        linear combination of the \emph{same} inputs: $+0.252 \pm 0.056$ on thin
        links, $+0.075 \pm 0.019$ at four smearing levels. An untrained network
        of identical architecture sits \emph{below} the classical GEVP at every
        point, and a width control shows $16 \to 24$ channels changes nothing.
        None of the three pre-registered readings survives its own test
        unmodified; the honest fourth statement is that the gap is largest where
        the classical basis is weakest and decays toward the saturation
        ceiling.
  \item \textbf{What the network found is a mediocre direction that is valuable
        in superposition, and it is mostly loop shapes.} The same operator lies
        $21.1\%$, $12.8\%$ and $2.5\%$ outside the $4$-, $7$- and
        $21$-dimensional classical spans; deleting the out-of-span part costs
        the entire advantage ($+0.097 \pm 0.021$ against the strong span, at
        unchanged mass); $80\%$ of it is rectangular Wilson loops the network
        rediscovered, leaving $\approx2.5\%$ of the operator's norm that no
        member of a $21$-operator basis expresses and that still carries
        $+0.021 \pm 0.009$. Alone the residual is a poor operator
        ($\Azero \approx 0.42$--$0.46$): it wins by constructive interference,
        carrying about twice the amplitude share its norm share would
        suggest.
  \item \textbf{The control decides which of those statistics means
        ``learned''.} An untrained network is $16$--$74\%$ outside the same span
        --- further out than the trained one --- so the norm fraction is
        architectural and carries no evidence of learning. Training moves the
        operator \emph{closer} to the classical span while making it better.
        What separates learned from architectural is the amplitude share
        ($Z_r/Z_G = 0.146$ against $0.047$, or $0.41$ against $0.10$ per unit of
        norm), the fact that the residual becomes loop-shaped at all ($80\%$
        against $0.3\%$ inside the richer span), and above all the sign of
        $\dAzero(\text{net}-P)$, which is $+0.097 \pm 0.021$ trained and
        $-0.091 \pm 0.042$ untrained: for a trained network, deleting the
        out-of-span content costs the advantage; for an untrained one, deleting
        it helps.
\end{itemize}

What remains is a ceiling rather than a missing measurement. Every strong arm on
these ensembles reads $\Azero \ge 0.90$ against a bound of $1$, so the
differences that separate the methods are being measured in the last tenth of
the available range, and no number of configurations changes that. The next real
step is a harder observable --- the first excited state, gated by a test that
costs minutes --- rather than a better estimator for this one.

\vfill
\footnotesize
\noindent\textbf{Reproducing the numbers.} All three studies are offline, CPU,
and read only files tracked in \texttt{dumps/}. The figures live under
\texttt{results/}, which is gitignored but regenerable in minutes:
Figure~\ref{fig:fairfight} from
\texttt{SFF\_TRUNCATE=1 SFF\_BASES=1 python scripts/su2\_fair\_fight.py <dump>},
Figure~\ref{fig:curve} from \texttt{python scripts/input\_architecture\_curve.py}
(which also writes Table~\ref{tab:curve} as a paste-ready \texttt{.tex}
fragment), and both panels of Figure~\ref{fig:decomp} from
\texttt{python scripts/operator\_decomposition.py} with
\texttt{-{}-basis=<obars>:deep}, \texttt{-{}-shape-span=<obars>:full} and, for
the untrained arm, \texttt{-{}-m-ref=<trained mass>}. The source of truth for
every number quoted here is \texttt{notes/fable5.1\_10-09\_audit.md}
(\S1--\S2 pre-registration, \S8.1 estimator and shape verdict, \S8.2 curve,
\S8.3 decomposition and controls) and
\texttt{notes/operator\_decomposition.md} (\S2 theory, \S4 published results,
\S5 controls, \S6.1 the withdrawn extrapolation, \S7 limitations).

\end{document}
